\documentclass[11pt]{article}
\usepackage[margin=1in]{geometry}
\usepackage{amsmath,amssymb,amsthm}
\usepackage{graphicx}
\usepackage[hidelinks]{hyperref}

\newtheorem{theorem}{Theorem}
\newtheorem{lemma}{Lemma}
\newtheorem{remark}{Remark}

\newcommand{\Baire}{\mathcal B_1}
\newcommand{\clp}{\operatorname{cls}_p}
\newcommand{\RR}{\mathbb R}
\newcommand{\NN}{\mathbb N}
\newcommand{\co}{\operatorname{co}}

\title{Separable Strongly Rosenthal Compacta and the Admissible Envelope Reduction}
\author{agent\_lane\_16}
\date{2026-06-16}

\begin{document}
\maketitle

\section*{Status}

This packet records a likely valid partial result for the question of
Glasner and Megrelishvili \cite{GM0803}.  It proves the full separable
subcase:

\begin{center}
\emph{Every separable strongly Rosenthal compactum is admissible.}
\end{center}

It also gives the following exact reformulation:

\begin{center}
\emph{\(K\) is admissible if and only if \(K\) is a compact subspace of a
separable strongly Rosenthal compactum.}
\end{center}

Consequently the full source question is equivalent, after this packet, to the
separable-envelope problem: does every strongly Rosenthal compactum embed into
a separable strongly Rosenthal compactum?  The nonseparable case remains open
here.

\section*{Source Question}

The source paper defines a compactum \(K\) to be strongly Rosenthal when it
embeds into \(\Baire(X)\) for a compact metrizable space \(X\), and admissible
when it is contained in the pointwise closure of a bounded set of continuous
functions whose pointwise closure consists of Baire-one functions.  The source
then asks whether the two classes coincide.

\begin{center}
\includegraphics[width=\textwidth]{figures/open_problem_crop.png}
\end{center}

\section*{Definitions and Source Tools}

Let \(X\) be compact metrizable.  A bounded family \(F\subset C(X)\) is a
Rosenthal family if \(\clp(F)\subset \Baire(X)\).  In the source paper this is
equivalent, by the Talagrand/Rosenthal criterion, to saying that \(F\) contains
no sequence equivalent to the unit vector basis of \(\ell_1\) \cite{GM0803}.

A compactum is admissible if it is homeomorphic to a compact subspace of
\(\clp(F)\) for some Rosenthal family \(F\subset C(X)\) on a compact
metrizable \(X\).

We use two results from Glasner--Megrelishvili \cite{GM0803}.

\begin{itemize}
  \item Theorem 6.3: if \(F\subset C(Y)\) is a Rosenthal family on compact
  metrizable \(Y\), then \(L=\clp(F)\) embeds into \(V^{**}\), with the
  weak-star topology, for a separable Rosenthal Banach space \(V\).  Moreover,
  for the trivial action version used here, the embedding extends the original
  functions in \(F\) and represents evaluation on \(Y\) by functionals in
  \(V^*\).
  \item Theorem 6.16: a compactum is admissible if and only if it is
  homeomorphic to a weak-star closed subset of \(V^{**}\) for some separable
  Rosenthal Banach space \(V\).
\end{itemize}

\section*{Proof Intuition}

The previous push produced an admissible cover.  Given a separable compact
\(K\subset\Baire(X)\), choose a dense sequence \((f_n)\) in \(K\), compactify
the graph
\[
  x\mapsto (x,f_0(x),f_1(x),\ldots),
\]
and let \(g_n\) be the coordinate functions on the compactified graph.  Then
\((g_n)\) is a Rosenthal family of continuous functions, and
\(L=\clp(\{g_n\})\) maps continuously onto \(K\) by restriction to the
original graph.

The missing step was not to find a continuous section of this quotient.  The
new idea is to quotient the Banach space representing \(L\).  The source
representation theorem embeds \(L\) into \(V^{**}\) for a separable Rosenthal
space \(V\).  Let \(N\subset V\) be the closed subspace of vectors whose
represented functions vanish on the original graph points.  Passing from
\(V\) to \(W=V/N\) collapses precisely those elements of \(L\) that have the
same restriction to the original graph.  Therefore the quotient bidual
\(W^{**}\) contains a weak-star compact copy of \(K\).  The source
second-dual characterization then implies that \(K\) is admissible.

\section*{Theorem}

\begin{theorem}
Every separable strongly Rosenthal compactum is admissible.
\end{theorem}

\begin{proof}
Since the property is topological, compose with a homeomorphism
\(\RR\to(-1,1)\) coordinatewise if necessary.  Thus it is enough to consider a
compact separable set
\[
  K\subset[-1,1]^X\cap \Baire(X),
\]
where \(X\) is compact metrizable and \(K\) has the pointwise topology.  Choose
a countable pointwise-dense sequence \((f_n)_{n\in\NN}\) in \(K\).

Define
\[
  e:X\to X\times[-1,1]^{\NN},
  \qquad
  e(x)=(x,f_0(x),f_1(x),\ldots),
\]
and let
\[
  Y=\overline{e(X)}\subset X\times[-1,1]^{\NN}.
\]
The space \(Y\) is compact metrizable.  Let \(g_n:Y\to[-1,1]\) be the
\(n\)-th coordinate function, so \(g_n\in C(Y)\) and
\[
  g_n(e(x))=f_n(x) \qquad (x\in X).
\]

First we recall the cover argument.  The family \(F=\{g_n:n\in\NN\}\) is a
Rosenthal family on \(Y\).  Indeed, if some subsequence \((g_{n_j})\) were
equivalent to the unit vector basis of \(\ell_1\), then by the
Talagrand/Rosenthal finite-independence criterion there would be real numbers
\(a<b\) such that for every finite disjoint \(P,Q\subset\NN\) the inequalities
\[
  g_{n_j}(y)<a \quad (j\in P),
  \qquad
  g_{n_j}(y)>b \quad (j\in Q)
\]
hold at some \(y\in Y\).  For fixed \(P,Q\), these strict inequalities define
a nonempty open subset of \(Y\).  Since \(e(X)\) is dense in \(Y\), the same
pattern is realized at some \(e(x)\).  Hence
\[
  f_{n_j}(x)<a \quad (j\in P),
  \qquad
  f_{n_j}(x)>b \quad (j\in Q).
\]
Thus \((f_{n_j})\) contains an \(\ell_1\)-sequence, contradicting the fact
that \(K\) is pointwise compact in \(\Baire(X)\): every sequence in \(K\) has a
pointwise convergent subsequence, while an \(\ell_1\)-sequence of bounded
functions has no pointwise convergent subsequence.  Therefore \(F\) is a
Rosenthal family.

Set
\[
  L=\clp(F)\subset[-1,1]^Y.
\]
Then \(L\) is admissible.  Define
\[
  \pi:L\to K,
  \qquad
  \pi(h)(x)=h(e(x)).
\]
This map is continuous for the pointwise topologies.  It maps \(g_n\) to
\(f_n\), and compactness gives \(\pi(L)=K\): if \(f\in K\), choose a net
\((f_{n_i})\) converging pointwise to \(f\); a subnet of \((g_{n_i})\)
converges pointwise on \(Y\) to some \(h\in L\), and then
\(\pi(h)=f\).

Now apply Theorem 6.3 of \cite{GM0803} to the Rosenthal family
\(F\subset C(Y)\), with the trivial group action.  We obtain a separable
Rosenthal Banach space \(V\), a weak-star homeomorphic embedding
\[
  \nu_0:L\hookrightarrow V^{**},
\]
and a weak-star continuous map \(\alpha:Y\to V^*\) such that
\[
  h(y)=\langle \nu_0(h),\alpha(y)\rangle
  \qquad (h\in L,\ y\in Y).
\]
Moreover, \(\nu_0(g_n)\in V\) for every \(n\).

Let
\[
  N=\bigcap_{x\in X}\ker \alpha(e(x)) \subset V,
\]
and let \(q:V\to W:=V/N\) be the quotient map.  Since quotients of separable
Rosenthal Banach spaces are again separable Rosenthal, \(W\) is a separable
Rosenthal space.

Identify \(W^*\) with the annihilator
\[
  N^\perp=\{\varphi\in V^*: \varphi|_N=0\}.
\]
By the elementary annihilator identity,
\[
  N^\perp=\overline{\operatorname{span}}\{\alpha(e(x)):x\in X\}
  \quad\text{inside }V^*.
\]
We claim that \(q^{**}\circ\nu_0\) collapses exactly the fibers of \(\pi\).
Let \(h_1,h_2\in L\).  Then
\[
\begin{aligned}
  \pi(h_1)=\pi(h_2)
  &\iff h_1(e(x))=h_2(e(x)) \quad \forall x\in X\\
  &\iff \langle \nu_0(h_1)-\nu_0(h_2),\alpha(e(x))\rangle=0
     \quad \forall x\in X\\
  &\iff \nu_0(h_1)-\nu_0(h_2)\text{ vanishes on }N^\perp\\
  &\iff q^{**}\nu_0(h_1)=q^{**}\nu_0(h_2).
\end{aligned}
\]
The third equivalence uses that elements of \(V^{**}\) are norm-continuous on
\(V^*\), so vanishing on the spanning set implies vanishing on its closed
linear span.

It follows that the continuous map \(q^{**}\circ\nu_0:L\to W^{**}\) factors
through the quotient map \(\pi:L\to K\).  Hence there is a unique continuous
map
\[
  \sigma:K\to W^{**}
\]
such that
\[
  \sigma(\pi(h))=q^{**}\nu_0(h) \qquad (h\in L).
\]
The fiber calculation above also shows that \(\sigma\) is injective.  Since
\(K\) is compact and \(W^{**}\) is Hausdorff in the weak-star topology,
\(\sigma\) is a homeomorphism of \(K\) onto the compact set \(\sigma(K)\).
The image \(\sigma(K)\) is weak-star compact, hence weak-star closed and
bounded, in \(W^{**}\).

Thus \(K\) is homeomorphic to a weak-star closed subset of the second dual of
the separable Rosenthal Banach space \(W\).  By Theorem 6.16 of
\cite{GM0803}, \(K\) is admissible.
\end{proof}

\begin{theorem}
For a compact space \(K\), the following are equivalent.
\begin{enumerate}
  \item \(K\) is admissible.
  \item \(K\) is homeomorphic to a compact subspace of a separable strongly
  Rosenthal compactum.
\end{enumerate}
\end{theorem}

\begin{proof}
If \(K\) is admissible, then \(K\subset\clp(F)\) for a Rosenthal family
\(F\subset C(X)\) on a compact metrizable space \(X\).  Since \(C(X)\) is
separable, choose a countable norm-dense subset \(D\subset F\).  Then
\(\clp(D)=\clp(F)\): the inclusion \(\clp(D)\subset\clp(F)\) is immediate,
while every \(f\in F\) is the uniform, hence pointwise, limit of a sequence
from \(D\).  Hence \(L=\clp(D)\) is a separable strongly Rosenthal compactum
containing \(K\).

Conversely, suppose that \(K\) is a compact subspace of a separable strongly
Rosenthal compactum \(S\).  By Theorem 1, \(S\) is admissible.  The class of
admissible compacta is hereditary for compact subspaces, so \(K\) is
admissible.
\end{proof}

\section*{Limitations}

This settles the separable strongly Rosenthal case and the source paper's
commented-out separable-envelope characterization.  It does not settle the full
source question because a strongly Rosenthal compactum need not be separable as
a topological space, nor is it known here that every strongly Rosenthal compactum
embeds into a separable strongly Rosenthal compactum.  The proof of Theorem 1
uses a countable dense sequence in \(K\) in an essential way to build the
metrizable graph compactification and countable Rosenthal family.

The full nonseparable question should therefore be read in the sharper form
given by Theorem 2.  A direct attack using all points of a nonseparable
\(K\subset\Baire(X)\) gives a graph compactification over
\(X\times[-1,1]^K\), but that space is generally nonmetrizable and the
corresponding coordinate family is uncountable.  The Banach-space quotient
argument then naturally produces a representation over a nonseparable
Rosenthal space, whereas admissibility requires a separable Rosenthal space.
Thus the remaining obstruction is not the quotient step; it is the absence, in
this pass, of a theorem producing a separable strongly Rosenthal envelope for
an arbitrary strongly Rosenthal compactum.

\section*{Novelty and Search Notes}

This packet is promoted as a partial result for the original question and as a
full result for the separable subcase.  The agent searched the local run
indexes for arXiv:0803.2320, strongly Rosenthal, admissible compacta, and the
previous cover packet.  Targeted web searches for the phrases
\emph{separable strongly Rosenthal admissible compactum},
\emph{every separable strongly Rosenthal admissible}, and close variants found
no exact later statement in the bounded search.  The source TeX itself
contains a commented-out attempted implication for a subspace of a separable
strongly Rosenthal compactum, stopping after the pointwise approximation step;
the quotient argument above appears to fill the separable strongly Rosenthal
case by routing through the admissible graph cover.

A final full-case literature/reduction pass checked nearby countable coding,
analytic-subspace, and strong-boundedness machinery for separable Rosenthal
compacta \cite{DodosCoding,ADK,DodosStrong}.  These tools are powerful
precisely in the separable or analytic-family setting; they did not supply, in
this pass, a separable envelope theorem for an arbitrary nonseparable strongly
Rosenthal compactum.

Novelty should still be treated as tentative until a human checks nearby
Rosenthal compacta folklore and the exact published version of the
Glasner--Megrelishvili theorem.

\section*{Verification Notes}

No computational verification is involved.  The proof has three key review
points:

\begin{itemize}
  \item the graph-coordinate family \((g_n)\) is Rosenthal because any
  finite-independence pattern in \(Y\) transfers back to the dense original
  graph \(e(X)\);
  \item the source Theorem 6.3 representation is used with the trivial action
  and preserves the original coordinate functions as vectors of \(V\);
  \item for
  \(N=\bigcap_x\ker\alpha(e(x))\), the annihilator
  \(N^\perp\) is exactly the closed linear span of the graph evaluations,
  so \(q^{**}\) identifies exactly the fibers of \(\pi:L\to K\).
\end{itemize}

\section*{Human Review Recommendation}

Treat this as a serious candidate proof of the separable version of the
Glasner--Megrelishvili question.  The most important line to audit is the
fiber calculation for \(q^{**}\circ\nu_0\).  If that calculation is accepted,
the rest is a direct application of the source paper's second-dual
characterization of admissibility.

\begin{thebibliography}{9}

\bibitem{GM0803}
E. Glasner and M. Megrelishvili,
\emph{Representations of dynamical systems on Banach spaces not containing
\(\ell_1\)},
arXiv:0803.2320.

\bibitem{BFT}
J. Bourgain, D. H. Fremlin and M. Talagrand,
\emph{Pointwise compact sets in Baire-measurable functions},
Amer. J. Math. 100 (1978), no. 4, 845--886.

\bibitem{Rosenthal}
H. P. Rosenthal,
\emph{A characterization of Banach spaces containing \(\ell^1\)},
Proc. Nat. Acad. Sci. USA 71 (1974), 2411--2413.

\bibitem{DFJP}
W. J. Davis, T. Figiel, W. B. Johnson and A. Pelczynski,
\emph{Factoring weakly compact operators},
J. Funct. Anal. 17 (1974), 311--327.

\bibitem{DodosCoding}
P. Dodos,
\emph{Codings of separable compact subsets of the first Baire class},
arXiv:0805.2026.

\bibitem{ADK}
S. A. Argyros, P. Dodos and V. Kanellopoulos,
\emph{A classification of separable Rosenthal compacta and its applications},
arXiv:0805.2032.

\bibitem{DodosStrong}
P. Dodos,
\emph{A strong boundedness result for separable Rosenthal compacta},
arXiv:0805.2034.

\end{thebibliography}

\end{document}
